% !Mode:: "TeX:UTF-8"

%%% 说明: 此部分需要自己填写的内容:  论文创新点.

%%% ``郑重申明'' 和``论文使用授权协议书'' 分别在 statement.tex 和 protocol.tex 中, 无需改动.

%%%%%%%%%%%%%%%%%%%%%%%%%%%%%
%%% -------------  英文封面 (无需改动)-------------   %%%
%%%%%%%%%%%%%%%%%%%%%%%%%%%%%
\thispagestyle{empty}
\renewcommand{\baselinestretch}{1.5}  %下文的行距
\vspace*{0.5cm}

\begin{center}{\zihao{2} \the\Etitle \par}\end{center}

\vfill

\begin{center}
\zihao{4}
\begin{tabular}{ r l }
 Candidate:      &  {\sc \the\Eauthor}      \\
 Supervisor:     &  {\sc \the\Esupervisor}   \\
 Major:          & \the\Emajor  \\
 Speciality:     & \the\Especiality
\end{tabular}

\vspace*{2cm}
\begin{center}
  \iflib % 向图书馆提交电子文档, 使用黑白校徽.
  \includegraphics[height=4cm]{whu.eps}       %%  黑白的. 很小, 只有 10k.
  \else
     \ifprint % 文档打印, 使用黑白校徽.
  \includegraphics[height=4cm]{whu.eps}       %%  黑白的.
  \else
  \includegraphics[height=4cm]{whulogo.eps} %%  彩色的.
  \fi
  \fi
\end{center}


\zihao{-2}
\the\Schoolname\\
{\sc Wuhan University}

\vspace*{1.0cm}

\the\Edate

\end{center}
%%%%%%%--判断是否需要空白页-----------------------------
  \iflib
  \else
  \newpage
  \thispagestyle{empty}
  \cleardoublepage
  \fi
%%%%%%%-------------------------------------------------
{\pagestyle{empty}
%%%--- 加入``郑重声明'' --- %%%%%%%%%%%%%%%%%
\newpage

\centerline{\heiti\zihao{-2}{深圳大学本科毕业论文（设计）诚信声明}}


\vskip 3cm 


\begin{spacing}{2.0}
	\zihao{4}
本人郑重声明：所呈交的毕业论文（设计），题目《融合内容信息的单类协同过滤推荐算法研究》是本人在指导教师的指导下，独立进行研究工作所取得的成果。对本文的研究做出重要贡献的个人和集体，均已在文中以明确方式注明。除此之外，本论文不包含任何其他个人或集体已经发表或撰写过的作品成果。本人完全意识到本声明的法律结果。

\vskip 3cm

{\flushright{
		毕业论文（设计）作者签名：\hspace{2.2cm}
		
		
		
		\hspace{7.5cm}{日期：\hspace{2cm}年 \hspace{.5cm}月\hspace{.5cm}日}\hspace{4cm} 
	}}
\end{spacing}


%%%%%%%%%%%%%%%%%
%%% ---加入``武汉大学学位论文使用授权协议书'' ---  %%%%%%
% !Mode:: "TeX:UTF-8"

%%%%%---武汉大学学位论文使用授权协议书---%%%%%%%%%%%%
%%%%%%%%%%%%%%%%%%%%%%%%%%%%%%%%%%%
%%%%%%%%%%%%%%%%%%%%%%%%%%%%%%%%%%%
\newpage\vspace*{20pt}\thispagestyle{empty}
\begin{center}{\zihao{-2}\heiti 武汉大学学位论文使用授权协议书}\end{center}
\par\vspace*{30pt}

本学位论文作者愿意遵守武汉大学关于保存、使用学位论文的管理办法及规定,
即：学校有权保存学位论文的印刷本和电子版, 并提供文献检索与阅览服务；
学校可以采用影印、缩印、数字化或其它复制手段保存论文；
在以教学与科研服务为目的前提下, 学校可以在校园网内公布部分及全部内容.
\begin{enumerate}[1、]
  \item  在本论文提交当年, 同意在校园网内提供前十六页浏览服务.
  \item  在本论文提交~$\Box$~当年/~$\Box$~一年/~$\Box$~两年
            /~$\Box$~三年, 同意在校园网内允许读者
            浏览并下载全文.(保密论文解密后遵守此规定)
\end{enumerate}

\vskip 15mm

论文作者(签名)：\raisebox{-1ex}{\underline{\makebox[5cm][c]{}}}
\vskip2em
				          				
学\qquad\qquad\quad 号：\raisebox{-1ex}{\underline{\makebox[5cm][c]{}}}
\vskip2em	
					
学\qquad\qquad\quad 院：\raisebox{-1ex}{\underline{\makebox[5cm][c]{}}}					

\vskip  2cm
\begin{flushright}
 日期：\hskip2cm 年\hskip1.2cm 月\hskip1.2cm 日
\end{flushright}

%%%%%%%%%%%%%%%%%%%%%%%%%%%%%%%%%%%%%%%
%%%%%%%--判断是否需要空白页-----------------------------
  \iflib
  \else
  \newpage
  \thispagestyle{empty}
  \cleardoublepage
  \fi
%%%%%%%-------------------------------------------------
%%%%%%%%%%%%%%%%%%
%%%%%%%%%%%%%%%%%%%%%%%%%%%%%%%
}




%%%%%%%%%%%%%%%%%%%%%%%%%%%%%%%
%%% ------------------- 论文创新点----------------------- %%%
%%%%%%%%%%%%%%%%%%%%%%%%%%%%%%%

\newpage\vspace*{20pt}\thispagestyle{empty}
\begin{center}{\zihao{-2}\heiti 论文创新点}\end{center}
\par\vspace*{30pt}
\baselineskip=20pt


  本文的主要创新点在于:

\begin{enumerate}[(1)]
  \item  给出了武汉大学博士毕业论文 \LaTeX{} 模板. 说明了编译方法.

  \item  强调了公式排版的一些细节. 并用例子说明.

  \item  强调了标点符号等的用法. 指出了初学者的一些常见错误.


\end{enumerate}

















%%%%%%%%%%%%%%%%%%%%%%%%%%%%%%%
%%%%%%%--请勿删除以下内容 -------------------------------%
%%%%%%%--判断是否需要空白页-----------------------------%
  \iflib
  \let\cleardoublepage\clearpage
  \else
  \newpage
  \thispagestyle{empty}
  \cleardoublepage
  \fi
%%%%%%%---------------------------------------------------%



